(1) 近轴薄透镜对光线角度的变化为 $\Delta\theta=-\Phi y$。两片紧贴时光线在两片之间几乎没有额外横向位移，角度变化线性相加，因此总光焦度 $\Phi=\Phi_1+\Phi_2$；若两片间距不可忽略，主平面移动后并非简单代数和。
(2) Abbe数定义将材料的主波长折射率与两谱线折射率差联系起来，在题设线性薄透镜模型中，一级轴向色差对应总光焦度差 $\Delta\Phi=\Phi_1/V_1+\Phi_2/V_2$。要求两目标谱线同焦，需 $\Phi_1/V_1+\Phi_2/V_2=0$，说明两种透镜的光焦度必须异号。
(3) 联立 $\Phi_1+\Phi_2=\Phi$ 与 $\Phi_1/V_1+\Phi_2/V_2=0$，解得 $\Phi_1=\Phi V_1/(V_1-V_2)$，$\Phi_2=-\Phi V_2/(V_1-V_2)$。代回两个约束，分别得到 $\Phi$ 和零一级色差，符合既定总焦距要求。
(4) 若 $V_1>V_2$ 且整体为会聚 $\Phi>0$，则 $\Phi_1>0$、$\Phi_2<0$，即低色散材料正透镜配高色散材料负透镜。该设计仅让指定两谱线焦点重合，中间波长仍可能存在二级色谱，真实厚透镜及非零间隔还需计入球差、主平面与高阶色散。
